% Developmental Cell: neutral adaptation of an inspected 2026 Technology article.
% All prose, authors and diagrams below are original placeholders.
\PassOptionsToPackage{no-math}{fontspec}
\documentclass[UTF8,fontset=none,twoside,scheme=plain]{ctexart}
\usepackage[paperwidth=602.986bp,paperheight=782.986bp,left=53.121bp,right=59.754bp,top=99bp,bottom=50.186bp,headheight=32bp,headsep=15bp,footskip=27.4bp]{geometry}
% Shared dependencies and neutral material. Compile with XeLaTeX.
% No publisher logos, paper text, data, figures, or download metadata are included.
\usepackage{amsmath,amssymb,graphicx,xcolor,tikz,fancyhdr,needspace,xurl,hyperref,balance,etoolbox}
% The legacy balance package omits the first split box's final baseline depth.
% Reserve that measured depth rather than hiding a CJK overfull-box warning.
\makeatletter
\patchcmd{\@BAlancecol}{\setbox\@leftcolumn=\vbox to \@tempdima}{\advance\@tempdima by \dp1\relax\setbox\@leftcolumn=\vbox to \@tempdima}{}{\PackageWarning{cns-template}{Balance depth patch unavailable; inspect final columns}}
\makeatother
\usetikzlibrary{arrows.meta,positioning}
\usepackage[protrusion=true,expansion=false]{microtype}
\IfFontExistsTF{Noto Serif CJK SC}{%
  \setCJKmainfont{Noto Serif CJK SC}%
}{%
  \setCJKmainfont[BoldFont=FandolSong-Bold.otf,ItalicFont=FandolKai-Regular.otf]{FandolSong-Regular.otf}%
}
\IfFontExistsTF{Noto Sans CJK SC}{%
  \setCJKsansfont{Noto Sans CJK SC}%
}{%
  \setCJKsansfont[BoldFont=FandolHei-Bold.otf]{FandolHei-Regular.otf}%
}
\setmainfont{TeX Gyre Termes}\setsansfont{TeX Gyre Heros}
\setmonofont{TeX Gyre Cursor}
\definecolor{JournalBlue}{HTML}{0097C9}
\definecolor{JournalRed}{HTML}{B83443}
\definecolor{ScienceRed}{HTML}{CD1021}
\definecolor{CiteBlue}{HTML}{0000EE}
\hypersetup{colorlinks=true,linkcolor=JournalBlue,urlcolor=JournalBlue,
  pdftitle={CNS论文汉化：中性版式示例},pdfauthor={示例模板；非期刊官方模板},
  pdfsubject={Neutral source-inspired Chinese layout demonstration; no real research claims}}
\urlstyle{same}\raggedbottom
\setlength{\parindent}{1em}\setlength{\parskip}{0pt}
\setlength{\emergencystretch}{2em}\tolerance=2200
\XeTeXlinebreaklocale "zh"\XeTeXlinebreakskip=0pt plus .1pt
\pagestyle{fancy}\fancyhf{}
\renewcommand{\headrulewidth}{0pt}\renewcommand{\footrulewidth}{0pt}
\newcommand{\PaperTitle}{中文题名占位：从组织结构到证据解释}
\newcommand{\SampleAuthors}{Author A\textsuperscript{1}, Author B\textsuperscript{2,*}}
\newcommand{\SampleAffiliations}{\textsuperscript{1}作者单位一（示例占位）。\quad
  \textsuperscript{2}作者单位二（示例占位）。\par
  *通信作者信息待填。DOI、收稿和出版日期均应从目标原文核实。}
\newcommand{\SampleAbstract}{本段是中性摘要占位文字，不陈述任何真实研究发现。它演示中文题名、作者信息、摘要和双栏正文之间的层级关系。实际使用时，应完整替换为目标论文的忠实中文译文，并保留研究对象、方法、结果和结论的限定条件。图注、参考文献、作者声明和来源信息均属于独立的覆盖项。模板中的图形仅说明排版结构，不表示生物学机制、实验结果或统计关系。}
\newcommand{\DemoCaption}{\textbf{图1. 中性证据链示意（占位图，无实验数据）。}
  \textbf{A}，输入节点表示需要核对的原文材料。
  \textbf{B}，处理节点表示逐段转换及复核。
  \textbf{C}，输出节点表示正文、图注和书目等独立交付部分。
  箭头仅表示本示例的阅读顺序，不代表因果关系、效应方向或数量变化。
  该图由模板自行绘制，所有节点均为文字占位；无样本量、误差线、统计检验或比例尺。
  实际论文图注应逐项保留原稿报告的样本量、单位、检验方法、限定及图形署名。}
\newcommand{\DemoVector}{%
 \resizebox{\linewidth}{!}{%
 \begin{tikzpicture}[font=\sffamily\small,node distance=18mm,
   box/.style={draw=black!70,rounded corners=2pt,fill=black!3,minimum width=26mm,minimum height=16mm,align=center},
   flow/.style={-{Stealth[length=2.5mm]},line width=.7pt}]
  \node[box] (a) {\textbf{A}\quad 输入\\原文材料};
  \node[box,right=of a] (b) {\textbf{B}\quad 处理\\转换与复核};
  \node[box,right=of b] (c) {\textbf{C}\quad 输出\\中文译稿};
  \draw[flow] (a)--(b);\draw[flow] (b)--(c);
  \node[below=11mm of a,align=center] {标题与摘要\\正文与图表};
  \node[below=11mm of b,align=center] {保留限定\\检查单位};
  \node[below=11mm of c,align=center] {可追溯引文\\明确原文信息};
  \node[below=29mm of b,font=\sffamily\footnotesize,text=black!65] {示意占位；不包含真实实验数据};
 \end{tikzpicture}}%
}
\newcommand{\SampleBibliography}{%
 \section*{参考文献（示例占位）}
 {\fontsize{7.6bp}{10bp}\selectfont\setlength{\parindent}{0pt}\interlinepenalty=10000
 \hypertarget{ref:1}{}1. Author A. (20XX). Example reference title. \textit{Journal placeholder}. 卷、页码与DOI待填。此条仅供演示引文跳转。\par\smallskip
 \hypertarget{ref:2}{}2. Author B. (20YY). Another example reference title. \textit{Journal placeholder}. 出版信息待填；不要把占位条目带入正式译稿。\par}
}
\newcommand{\SampleSourceInfo}{%
 \section*{原文信息}
 以下字段全部为示例占位：作者；英文原标题；期刊；年份；卷期；页码或文章号；DOI或来源链接。实际使用时必须逐项替换并核对。\par
 本文件是\JournalName 风格的非官方中文译稿样板，依据一次已核验文章版式进行结构化改编，不代表该期刊所有年代或文体的统一标准，也不保证像素级复刻。无真实论文内容或期刊标识图像。\par
}
% Deliberately neutral demo prose, not translated from a published paper.
\newcommand{\SampleOpening}{%
 \section*{引言：示例正文与引文}
 这是一段用于检查双栏阅读节奏的中性示例文字。段落中的术语、缩写和标点需要在实际项目中与源文保持一致。中文译稿应呈现原论文的完整论述，而不是把多段证据压缩成一段摘要。这里的引文只链接到本文件末尾明确标为占位的参考文献\SampleCiteOne。\par
 若原句仅表示相关、可能或提示，译文也应保留相同的证据强度。数字、单位、正负号、上下标和范围端点应逐项核对。原文存在内部矛盾时，应分别保留相应位置的陈述，再用简短译注说明；不能根据常识把不同的数字静默统一。\par
 \subsection*{内容层级与跨页衔接}
 此小标题演示二级层级。实际论文可能采用不同的小标题密度，模板不会替代对原文结构的识别。双栏文章需按左栏、右栏和跨页续句的顺序读取。若段落在页末中断，应在完成整句核对后接入下一页，避免把图中文字误当正文，或重复翻译同一段。\par
 图表具有独立的信息范围。正文引用到某幅图，并不意味着其图注可以省略。图注中的每个分面、比较对象、时间点、样本量及统计说明都应对应源文。英文图内标签可以保留，并在图旁或文末提供中文释义；实际采用哪种安排，需要依据图形尺寸和译文长度判断\SampleCiteTwo。\par
 \subsection*{来源与可维护数据}
 本段继续演示正文密度。可编辑文本和图注数据应与最终版面分开维护，以便修订时准确定位变化。若只调整版式，应能证明正文、图注和书目内容没有被意外改写。更换字体或纸张后，须重新检查断行、图注与图片的相对位置以及页脚可读性。\par
 文内引用与书目之间应有明确对应关系。这里的两条演示书目没有真实DOI，也没有虚构的实验信息。使用模板时可增删条目，但必须同时更新编号、链接目标和引用位置。页码只表示中文译稿自己的分页，不用于冒充原出版页码。\par
 \subsection*{结果叙述的演示占位}
 本节的文字仅用于呈现长段落在双栏页面中的流动，不报告真实结果。实际翻译应保留观察对象、实验条件和结局指标之间的关系。例如，原文分别描述体外实验与动物实验时，不能把两类结果合并为同一层面的结论；研究使用的对照、时间顺序和分析范围也必须保留。图表中的比较范围不应扩大为对整个人群或全部疾病阶段的判断。\par
 论述中的转折与限制同样属于正文。若研究先提出一种解释，再讨论另一种可能，译文应沿用这一推进顺序。作者明确写出的不确定性应能在中文中被识别；如果尚未确定某个现象的方向或机制，就不能改成已经建立的因果关系。对于术语近似但研究层级不同的概念，可在首次出现时保留缩写，后文保持统一。\par
 \subsection*{方法信息与可复核性}
 方法部分往往包含密集数字和专有名词。翻译时须区分浓度、总量、频率和处理时长，不能因为中文句式调整而改变参数的归属。试剂名称、软件版本和数据访问标识通常保留原样；上下限、排除规则和显著性标准则需要逐字符核对。这里没有提供任何真实实验参数，避免把模板占位误认为可以直接执行的方法。\par
 表格与补充材料的引用也需要保留原层级。正文对补充图的指向并不意味着模板中已经附有该补充图；若源文件只包含主文，应如实标注所交付的范围。图表编号、附录顺序和参考文献数量都应由目标文章决定，本示例中的结构不构成固定配额，也不要求所有论文采用同样的章节。\par
}
\newcommand{\SampleClosing}{%
 \section*{讨论：版式适配与范围}
 上方矢量图是模板原创的排版占位图，不来自任何论文。图形与完整演示图注被放在同一页，便于读者直接核对分面说明。真实项目应使用经过授权或用户提供的原图素材，并保留图注署名与许可信息。尺寸不足时，可适度缩放或调整栏宽，不以删除图注解决空间问题。\par
 模板仅记录一次已核验文章的可观察设计特征。即使文章出自同一期刊，不同年代、研究论文与综述、校正清样与正式版本也可能采用不同的页眉和引文形式。应以用户提供的实际PDF为准，不能仅凭期刊名称推断配色、字号或引用方式\SampleCiteOne。\par
 \subsection*{演示检查项}
 本样板用于展示可编译性、中文字形、标题层级、跨栏图注、页眉页脚和内部链接。图中没有实际数据，因此不应被当作科学插图引用。替换内容后，还需检查所有页的渲染结果，确认无截字、碰撞、异常空白和遗漏段落。\par
 \SampleBibliography
 \SampleSourceInfo
}


\usepackage{multicol,longtable,booktabs,array}
\newcommand{\JournalName}{Developmental Cell}
\renewcommand{\familydefault}{\sfdefault}
\definecolor{DevCellBlue}{HTML}{0097D1}
\definecolor{DevCellRed}{HTML}{AB4D4C}
\setlength{\columnsep}{13.2bp}
\renewcommand{\normalsize}{\fontsize{9.1bp}{12.6bp}\selectfont}
\ctexset{section={format=\sffamily\bfseries\color{DevCellRed}\fontsize{9.4bp}{12.6bp}\selectfont,beforeskip=8bp,afterskip=3bp},subsection={format=\sffamily\bfseries\color{DevCellRed}\fontsize{9.1bp}{12.6bp}\selectfont,beforeskip=6bp,afterskip=2bp}}
\hypersetup{linkcolor=DevCellBlue,urlcolor=DevCellBlue,pdftitle={Developmental Cell风格：原创中性中文排版示例},pdfauthor={原创中性模板；非期刊官方模板},pdfsubject={Neutral source-inspired layout; no source-paper content or assets}}
\fancyhead[LO,RE]{\sffamily\color{DevCellBlue}\bfseries\fontsize{20bp}{25bp}\selectfont Developmental Cell}
\fancyhead[LE,RO]{\sffamily\color{black!35}\bfseries\fontsize{14bp}{18bp}\selectfont Technology}
\fancyfoot[LO,RE]{\sffamily\fontsize{6.8bp}{9bp}\selectfont 非官方中文译稿样板\quad 仅供学术交流}
\fancyfoot[LE,RO]{\sffamily\fontsize{7.5bp}{9bp}\selectfont Developmental Cell\quad 中文译页\ \thepage}
\newcommand{\DevCellTitle}{从材料核对到完整呈现：\\记录之间的对应关系}
\newcommand{\DevCellCite}[1]{\textsuperscript{\hyperlink{devcell-ref:#1}{#1}}}
\newcommand{\DevCellLabel}[1]{{\color{DevCellBlue}\fontsize{12bp}{16bp}\selectfont #1}\par\vspace{5bp}}
\newcommand{\DevCellBullet}[1]{\noindent\textcolor{DevCellBlue}{\textbullet}\quad #1\par\vspace{7bp}}
\newcommand{\DevCellGraphical}{%
 \resizebox{\linewidth}{!}{\begin{tikzpicture}[
  font=\sffamily\fontsize{11bp}{15bp}\selectfont,
  box/.style={draw=DevCellBlue!75,fill=DevCellBlue!6,rounded corners=2bp,minimum width=240bp,minimum height=65bp,align=center},
  arr/.style={-{Stealth[length=2.5mm]},line width=.9bp}]
  \node[box] (a) {\textbf{A\quad 来源材料}\\题名、正文、方法与完整图表};
  \node[box,below=22bp of a] (b) {\textbf{B\quad 逐项核对}\\数值、限定、对应关系与声明};
  \node[box,below=22bp of b] (c) {\textbf{C\quad 可追溯呈现}\\中文译稿、书目与原文信息};
  \draw[arr] (a)--(b);\draw[arr] (b)--(c);
 \end{tikzpicture}}%
}
\fancypagestyle{devcell-opening}{\fancyhead{}\fancyhead[L]{\sffamily\color{DevCellBlue}\bfseries\fontsize{20bp}{25bp}\selectfont Developmental Cell}}
\begin{document}\normalsize
% Page 1: separate graphical abstract, highlights and right-hand information.
{\noindent\bfseries\fontsize{22bp}{27bp}\selectfont\DevCellTitle\par}\vspace{13bp}
\noindent\begin{minipage}[t]{.59\textwidth}\vspace{0pt}
\DevCellLabel{Graphical abstract\quad 图形摘要}
\setlength{\fboxsep}{8bp}\setlength{\fboxrule}{.5bp}
\noindent\fbox{\begin{minipage}{\dimexpr\linewidth-17bp\relax}\DevCellGraphical\end{minipage}}\par\vspace{13bp}
\DevCellLabel{Highlights\quad 研究要点}
\DevCellBullet{正文、图表和方法分别清点，保留独立覆盖项。}
\DevCellBullet{比较方向、单位与不确定性按来源核对。}
\DevCellBullet{中文自然重排，完整图注与原图对应。}
\DevCellBullet{本页仅为原创中性示意，不报告研究结果。}
\end{minipage}\hfill
\begin{minipage}[t]{.385\textwidth}\vspace{0pt}
\DevCellLabel{Authors\quad 作者}
Author Placeholder A,\\Author Placeholder B\\（原创虚构作者占位）\par\vspace{17bp}
\DevCellLabel{Correspondence\quad 通信信息}
按目标来源填写通信作者。当前示例没有真实邮箱或联系方式。\par\vspace{17bp}
\DevCellLabel{In brief\quad 简介}
本例展示独立图形摘要页的左图与研究要点、右作者与简介结构。示意图中的节点仅表示本例的材料核对流程，不代表实验对象、数据、效应或因果关系。\par\vspace{17bp}
{\fontsize{7.6bp}{10bp}\selectfont 图形摘要、研究要点与正文摘要各自保留，不因文字相近而合并。正式稿的出版字段逐项从来源填写。}
\end{minipage}\par
\vfill
{\fontsize{7.6bp}{10bp}\selectfont 独立图形摘要页不替代正文摘要。出版身份以当前提供的原文件为依据。\par
仅供学术交流、出版物版权归原作者及出版商所有。}
\clearpage
% Page 2: Technology opening, continuous full-width summary and same-page body.
\thispagestyle{devcell-opening}
{\noindent\color{black!35}\bfseries\fontsize{14bp}{18bp}\selectfont Technology\par}\vspace{3bp}
{\noindent\bfseries\fontsize{22bp}{27bp}\selectfont\DevCellTitle\par}\vspace{10bp}
{\noindent\bfseries\fontsize{9bp}{12bp}\selectfont Author Placeholder A\textsuperscript{1,*}, Author Placeholder B\textsuperscript{2}\par}\vspace{4bp}
{\noindent\fontsize{7.5bp}{10.4bp}\selectfont\textsuperscript{1}原创中性作者单位一；\textsuperscript{2}原创中性作者单位二。\par *通信作者、DOI与出版日期按目标来源填写；此处无真实出版物对应信息。\par}\vspace{17bp}
{\noindent\color{DevCellRed}\bfseries SUMMARY\quad 摘要\par}\vspace{9bp}
{\fontsize{10bp}{14bp}\selectfont\color{DevCellBlue}
本文以原创中性材料演示完整记录的整理与核对。来源材料可以包含多种独立内容，整理时应先区分正文、图表、方法和声明，再恢复它们的阅读顺序。示意流程仅表达材料之间的对应关系，不表示真实实验、科学机制、效应方向或定量结果。\par
本例采用通栏题名、虚构作者单位及连续蓝色摘要，后接无衬线双栏正文。独立图形摘要、研究要点及简介在前页分别展示。这里不包含结构化摘要标签或独立背景意义框；正式稿应根据当前来源保留其实际结构。中文自然重排可以改变页数，数值、限定条件与引文关系仍须逐项核对。\par}
\vspace{15bp}
\begin{multicols}{2}
\section*{引言（原创中性示例）}
资料整理从区分独立记录开始。即使图形摘要、简介和正文摘要含有相近表述，它们也可能服务于不同的阅读顺序，应按原文分别呈现\DevCellCite{1}。

题名、作者、单位和日期共同确定来源版本。文件名、文件生成日期和正式出版日期可能不同，应以可见页面的明确出版信息为依据。
\section*{结果（版式示例）}
本例的结果段落只展示红色层级、无衬线双栏及蓝色数字上标引文\DevCellCite{2}。文中没有真实样本、统计检验、研究结论或医学建议。

完整呈现还包括方法细节、资源标识符、图注尾句及出版声明。这些内容不会因正文重排而省略；实际提供的补充材料须另行清点。
\end{multicols}
\vfill
{\fontsize{7.6bp}{10bp}\selectfont 本页全部为原创中性文字。普通刊名替代官方字标；中文字体与分页为有意适配。\par
仅供学术交流、出版物版权归原作者及出版商所有。}
\clearpage
% Page 3: full-width figure, complete caption and resource/method detail.
\newgeometry{left=59.754bp,right=53.1bp,top=99bp,bottom=50.186bp,headheight=32bp,headsep=15bp,footskip=27.4bp}
\section*{图文与方法（原创中性示例）}
\noindent\begin{minipage}{\textwidth}
\resizebox{\textwidth}{!}{\begin{tikzpicture}[
 font=\sffamily\fontsize{11bp}{15bp}\selectfont,
 box/.style={draw=DevCellBlue!75,fill=DevCellBlue!5,minimum width=125bp,minimum height=42bp,align=center},
 arr/.style={-{Stealth[length=2.5mm]},line width=.8bp}]
 \node[box] (a) {\textbf{A\quad 来源记录}\\正文、图表与方法};
 \node[box,right=28bp of a] (b) {\textbf{B\quad 对应核对}\\段落、数值与限定};
 \node[box,right=28bp of b] (c) {\textbf{C\quad 完整呈现}\\译文、资源与出处};
 \draw[arr] (a)--(b);\draw[arr] (b)--(c);
\end{tikzpicture}}\par\vspace{5bp}
{\fontsize{7.6bp}{9.2bp}\selectfont\textcolor{DevCellRed}{\bfseries 图1。记录对应流程（原创中性示意）。}
\textbf{A}，来源记录包含正文、图表与方法；\textbf{B}，核对段落、数值及其限定条件；\textbf{C}，呈现译文、资源和来源信息。箭头仅表示本例的阅读顺序，不表示因果、效应、数量关系或细胞分化。无真实数据、样本量、统计检验、误差线或比例尺。实际图注须完整保留全部面板、定义、单位及署名。\par}
\end{minipage}\vspace{6bp}
\section*{STAR方法：关键资源表（原创占位）}
{\fontsize{7.6bp}{9.7bp}\selectfont\setlength{\tabcolsep}{4bp}
\begin{longtable}{@{}p{.35\textwidth}p{.28\textwidth}p{\dimexpr.37\textwidth-16bp\relax}@{}}
\toprule 试剂或资源 & 来源 & 标识符\\\midrule\endfirsthead
\toprule 试剂或资源（续） & 来源 & 标识符\\\midrule\endhead
\bottomrule\endfoot
材料记录（原创占位） & 来源记录待填 & 货号／标识符待填\\
分析工具（原创占位） & 软件来源待填 & 版本／资源标识待填\\
数据记录（原创占位） & 数据存储位置待填 & 数据集编号待填\\
\end{longtable}}
\begin{multicols}{2}
\subsection*{方法与统计（结构示例）}
正式稿应分别保留实验模型、研究参与者详情、方法细节以及定量和统计分析。关键资源表逐行保留来源、货号及资源标识；长表采用可续页结构并重复列头。本例三行均为原创占位，不含真实试剂或实验参数。
\subsection*{补充内容与声明}
主文、STAR方法、补充图注、数据表与封面各自核对。源PDF中的原刊页码与物理页分别记录；被引用但未提供的独立附件不能视为已翻译。资助、作者贡献、利益冲突、数据和代码可用性及许可声明不得遗漏。
\section*{参考文献（原创占位）}
{\fontsize{7.6bp}{9.8bp}\selectfont\interlinepenalty=10000
\noindent\hypertarget{devcell-ref:1}{1.} Author Placeholder A. Complete record inventory. Original neutral reference example; no publication or DOI assigned.\par\smallskip
\noindent\hypertarget{devcell-ref:2}{2.} Author Placeholder B. Source and layout comparison. Original neutral reference example; no publication or DOI assigned.\par}
\section*{原文信息（占位字段）}
正式稿应按所给版本填写作者、英文原题、期刊、年份、卷期、页码或文章号、DOI及明确提供的日期。此处没有对应的真实论文，字段均待替换，不继承旧任务的出版身份。\par
仅供学术交流、出版物版权归原作者及出版商所有。
\end{multicols}
\end{document}
